At $\phi=0.376$--$0.393$, the same-exposure blue-minus-core residual is
$1176/852$ ppm in the primary/alternate pixels. Its conditional standard
deviation under the current fixed GP family is $355/345$ ppm, retaining
blue/core covariance. The fixed-planet empirical time-block sensitivity
instead spans $544$--$1052$ and $386$--$901$ ppm. Those large block estimates
include the possible physical residual pattern itself, so they are not proof
that the feature is noise. They demonstrate that these short segments cannot
independently calibrate a robust significance. Some variance estimates are
dominated by approximately one effective block.

No individual fine-bin B/C/R residual, or same-band difference between the
fine primary-product phase bins, keeps its sign in every displayed
two-standard-deviation block-proxy interval. In particular, the red residual
in the candidate interval is only $-216/-223$ ppm with conditional errors
$213/210$ ppm. A separate required red suppression is not established.
Residual ranges across 64 saved noise draws, including separately reprofiled
RV/amplitude, and changes in the local-refit covariance are smaller than these
block sensitivities. The local-refit proxy intervals retain the fixed baseline
center; they are a covariance sensitivity, not a newly optimized phase profile.
These controls do not resolve the uncertainty about omitted time/spectral
systematics.

The native phase-band correction fit has 17 observed columns and full rank
in its chosen rectangular basis; the last blue column is unobserved and
excluded. The complete-band displays retain only 12 or 15 independent rows,
depending on whether common or band-specific exposure support is required.
They cannot reconstruct all 17 corrections. For example, a positive blue
perturbation in one phase bin can cause a negative core response in the next
after the shared reference and GP are fitted. The finite recovery tests pass,
but they do not make the physical interpretation unique.
